% GENERATED FILE. Edit canonical lessons, then run npm run generate:books.
% canonical-source-hash: fnv1a64:4bec452805762f0a
% canonical-lessons: KA-C117-suttige, KA-C117-kattari, KA-C117-bacanige, KA-C117-sabunu, KA-C117-tavel

\chapter{Hammer, Scissors, Comb, Soap, Towel}
\label{ch:ka-hammer-scissors-comb-soap-towel}
\hlchaptermodality{117}

\begin{chapteropening}
\textbf{By the end of this chapter:} \emph{I can say a hammer, scissors, a comb, soap and a towel.}

Complete the last lesson of chapter 117: i can say a hammer, scissors, a comb, soap and a towel.
\end{chapteropening}

\section[\texorpdfstring{\kn{ಸುತ್ತಿಗೆ} (suttige)}{ಸುತ್ತಿಗೆ (suttige)}]{\kn{ಸುತ್ತಿಗೆ} (suttige) — a hammer}

\label{lesson:KA-C117-suttige}



\begin{quote}
Before the new one: say the Kannada for a lock, then the Kannada for a nail (to hammer).
\end{quote}

\subsection*{You'll want to know: \kn{ಸುತ್ತಿಗೆ}}
\textbf{\kn{ಸುತ್ತಿಗೆ}} — \emph{suttige} — \textquotedblleft{}a hammer\textquotedblright{}.

\begin{grammarlens}[title={What you've built}]
One more word for this chapter.
\end{grammarlens}

\subsection*{Your turn}
\begin{itemize}
\raggedright
  \item \emph{Say it:} \emph{suttige}
  \item \emph{Say it:} \emph{suttige}, once more
  \item \emph{Say it:} say \emph{suttige}
  \item \emph{Recall:} say the Kannada for a lock, then the Kannada for a nail (to hammer), then say \emph{suttige} again
\end{itemize}

\begin{tcolorbox}[breakable,colback=teal!4,colframe=teal!35!black,title={Before you move on}]
What does \kn{ಸುತ್ತಿಗೆ} mean? (\textquotedblleft{}A hammer\textquotedblright{}.) Say it once more.
\end{tcolorbox}
\section[\texorpdfstring{\kn{ಕತ್ತರಿ} (kattari)}{ಕತ್ತರಿ (kattari)}]{\kn{ಕತ್ತರಿ} (kattari) — scissors}

\label{lesson:KA-C117-kattari}



\begin{quote}
Before the new one: say the Kannada for a nail (to hammer), then the Kannada for a hammer.
\end{quote}

\subsection*{You'll want to know: \kn{ಕತ್ತರಿ}}
\textbf{\kn{ಕತ್ತರಿ}} — \emph{kattari} — \textquotedblleft{}scissors\textquotedblright{}.

\begin{grammarlens}[title={What you've built}]
One more word for this chapter.
\end{grammarlens}

\subsection*{Your turn}
\begin{itemize}
\raggedright
  \item \emph{Say it:} \emph{kattari}
  \item \emph{Say it:} \emph{kattari}, once more
  \item \emph{Say it:} say \emph{kattari}
  \item \emph{Recall:} say the Kannada for a nail (to hammer), then the Kannada for a hammer, then say \emph{kattari} again
\end{itemize}

\begin{tcolorbox}[breakable,colback=teal!4,colframe=teal!35!black,title={Before you move on}]
What does \kn{ಕತ್ತರಿ} mean? (\textquotedblleft{}Scissors\textquotedblright{}.) Say it once more.
\end{tcolorbox}
\section[\texorpdfstring{\kn{ಬಾಚಣಿಗೆ} (bācaṇige)}{ಬಾಚಣಿಗೆ (bācaṇige)}]{\kn{ಬಾಚಣಿಗೆ} (bācaṇige) — a comb}

\label{lesson:KA-C117-bacanige}



\begin{quote}
Before the new one: say the Kannada for a hammer, then the Kannada for scissors.
\end{quote}

\subsection*{You'll want to know: \kn{ಬಾಚಣಿಗೆ}}
\textbf{\kn{ಬಾಚಣಿಗೆ}} — \emph{bācaṇige} — \textquotedblleft{}a comb\textquotedblright{}.

\begin{grammarlens}[title={What you've built}]
One more word for this chapter.
\end{grammarlens}

\subsection*{Your turn}
\begin{itemize}
\raggedright
  \item \emph{Say it:} \emph{bācaṇige}
  \item \emph{Say it:} \emph{bācaṇige}, once more
  \item \emph{Say it:} say \emph{bācaṇige}
  \item \emph{Recall:} say the Kannada for a hammer, then the Kannada for scissors, then say \emph{bācaṇige} again
\end{itemize}

\begin{tcolorbox}[breakable,colback=teal!4,colframe=teal!35!black,title={Before you move on}]
What does \kn{ಬಾಚಣಿಗೆ} mean? (\textquotedblleft{}A comb\textquotedblright{}.) Say it once more.
\end{tcolorbox}
\section[\texorpdfstring{\kn{ಸಾಬೂನು} (sābūnu)}{ಸಾಬೂನು (sābūnu)}]{\kn{ಸಾಬೂನು} (sābūnu) — soap}

\label{lesson:KA-C117-sabunu}



\begin{quote}
Before the new one: say the Kannada for scissors, then the Kannada for a comb.
\end{quote}

\subsection*{You'll want to know: \kn{ಸಾಬೂನು}}
\textbf{\kn{ಸಾಬೂನು}} — \emph{sābūnu} — \textquotedblleft{}soap\textquotedblright{}.

\begin{grammarlens}[title={What you've built}]
One more word for this chapter.
\end{grammarlens}

\subsection*{Your turn}
\begin{itemize}
\raggedright
  \item \emph{Say it:} \emph{sābūnu}
  \item \emph{Say it:} \emph{sābūnu}, once more
  \item \emph{Say it:} say \emph{sābūnu}
  \item \emph{Recall:} say the Kannada for scissors, then the Kannada for a comb, then say \emph{sābūnu} again
\end{itemize}

\begin{tcolorbox}[breakable,colback=teal!4,colframe=teal!35!black,title={Before you move on}]
What does \kn{ಸಾಬೂನು} mean? (\textquotedblleft{}Soap\textquotedblright{}.) Say it once more.
\end{tcolorbox}
\section[\texorpdfstring{\kn{ಟವೆಲ್} (ṭavel)}{ಟವೆಲ್ (ṭavel)}]{\kn{ಟವೆಲ್} (ṭavel) — a towel}

\label{lesson:KA-C117-tavel}



\begin{quote}
Before the new one: say the Kannada for a comb, then the Kannada for soap.
\end{quote}

\subsection*{You'll want to know: \kn{ಟವೆಲ್}}
\textbf{\kn{ಟವೆಲ್}} — \emph{ṭavel} — \textquotedblleft{}a towel\textquotedblright{}.

\begin{grammarlens}[title={What you've built}]
One more word for this chapter.
\end{grammarlens}

\subsection*{Your turn}
\begin{itemize}
\raggedright
  \item \emph{Say it:} \emph{ṭavel}
  \item \emph{Say it:} \emph{ṭavel}, once more
  \item \emph{Say it:} say \emph{ṭavel}
  \item \emph{Recall:} say the Kannada for a comb, then the Kannada for soap, then say \emph{ṭavel} again
\end{itemize}

\begin{tcolorbox}[breakable,colback=teal!4,colframe=teal!35!black,title={Before you move on}]
What does \kn{ಟವೆಲ್} mean? (\textquotedblleft{}A towel\textquotedblright{}.) Say it once more.
\end{tcolorbox}
